\section{A probability law from transfer growth}
\label{sec:weighted-law}

The disk estimate in Proposition~\ref{prop:disk-transfer} reduces our task to
proving that $\rho\leq 2$.  We now suppose that $\rho>2$ and construct a
probability measure on the normalized family $\KK$.  Its covariance under
affine changes of coordinates will turn geometric quantities into expectations
at the single point $i$.

We use the operators $A_z$ and $L$ from the preceding section.  Measures act
on their left: if $\mu$ is a finite measure, then
$(\mu A_z)(\Phi)=\mu(A_z\Phi)$ for $\Phi\in C(\KK)$.  Every positive bounded
functional on $C(\KK)$ is identified with its finite Borel measure by the
Riesz representation theorem \cite[Theorem~7.2]{Folland1999}.

\begin{lemma}[Weighted affine covariance]
\label{lem:weighted-law}
Suppose that $\rho>2$, and put $\beta=\log_2\rho>1$.  There is a Borel
probability measure $P$ on $\KK$ such that, for every $z=x+iy\in\HH$ and
every nonnegative Borel function $\Phi:\KK\to[0,\infty]$,
\begin{equation}
 \EE\!\left[|Q_F(z)|^2\Phi(T_zF)\right]
   =y^{-\beta}\EE\Phi(F),
 \label{eq:weighted-covariance}
\end{equation}
where $\EE$ denotes expectation with respect to $P$.
\end{lemma}

\begin{proof}
We first obtain a probability eigenmeasure for $L$.  The dyadic formula for
$L^n$ will then give one horizontal period.  Averaging over that period and
over one logarithmic scale interval produces
\eqref{eq:weighted-covariance}.

\paragraph{The eigenmeasure.}
For $\tau>\rho$, the root limit defining $\rho$ implies that
\[
 R_\tau=\sum_{n=0}^{\infty}\tau^{-n}L^n
\]
converges in operator norm.  Set
\[
 g_\tau=R_\tau\mathbf 1,
 \qquad M_\tau=\max_{F\in\KK}g_\tau(F).
\]
Positivity gives $1\leq g_\tau\leq M_\tau$, and the maximum is attained
because $g_\tau$ is continuous and $\KK$ is compact.  We claim that
\begin{equation}
 M_\tau\longrightarrow\infty\qquad(\tau\downarrow\rho).
 \label{eq:resolvent-divergence}
\end{equation}
Indeed,
\[
 Lg_\tau=\tau(g_\tau-\mathbf 1)
 \leq\tau(1-M_\tau^{-1})g_\tau.
\]
Iteration, using positivity and $\mathbf 1\leq g_\tau$, gives
\[
 L^n\mathbf 1
 \leq\bigl[\tau(1-M_\tau^{-1})\bigr]^n g_\tau.
\]
Since $\|L^n\|=\|L^n\mathbf 1\|_\infty$ by
Lemma~\ref{lem:transfer-growth}, it follows that
\[
 \|L^n\|\leq M_\tau
      \bigl[\tau(1-M_\tau^{-1})\bigr]^n.
\]
If $M_{\tau_j}\leq M<\infty$ along a sequence $\tau_j\downarrow\rho$,
taking the $n$th-root limit for each fixed $j$ would give
\[
 \rho\leq\tau_j(1-M_{\tau_j}^{-1})
       \leq\tau_j(1-M^{-1}).
\]
Letting $j\to\infty$ contradicts $\rho>0$, proving
\eqref{eq:resolvent-divergence}.

Choose a maximizer $F_\tau$ and define
\[
 \mu_\tau=\frac{\delta_{F_\tau}R_\tau}{M_\tau}.
\]
This is a positive functional of mass one, hence a Borel probability measure.
The norm-convergent series gives $R_\tau L=\tau(R_\tau-I)$, where $I$ is
the identity operator, so
\[
 \mu_\tau L=\tau\mu_\tau-
                \frac{\tau}{M_\tau}\delta_{F_\tau}.
\]
Choose a sequence $\tau\downarrow\rho$.  Since $\KK$ is compact metric,
$C(\KK)$ has a countable uniformly dense subset.  Diagonal selection gives
a subsequence along which the integrals of those functions converge.
The bound $\|\mu_\tau\|=1$ extends convergence to every function in
$C(\KK)$.  The limiting functional is positive and has mass one, hence is
a probability measure $\mu$ by Riesz representation.  For each
$\Phi\in C(\KK)$, the function $L\Phi$ is continuous, and the final term
above has absolute value at most $\tau\|\Phi\|_\infty/M_\tau$.
Consequently
\begin{equation}
 \mu L=\rho\mu.
 \label{eq:transfer-eigenmeasure}
\end{equation}

\paragraph{A horizontal period.}
For $v=2^{-n}$ and $u_j=-1+(2j+1)v$, the formula from
Lemma~\ref{lem:transfer-growth} reads
\[
 L^n=v\sum_{j=0}^{2^n-1}A_{u_j+iv}.
\]
Right composition with $A_{2+i}$ replaces $u_j$ by $u_j+2v$.
Thus \eqref{eq:transfer-eigenmeasure} gives the exact cancellation
\begin{equation}
 \mu A_{2+i}-\mu
 =\rho^{-n}v\mu
    \bigl(A_{1+v+iv}-A_{-1+v+iv}\bigr).
 \label{eq:horizontal-endpoints}
\end{equation}
Both remaining abscissae have absolute value at most $2$.  The bound
$|Q_F(u+iv)|\leq C/v$ from Lemma~\ref{lem:compact-class} therefore implies
\[
 \|\mu A_{2+i}-\mu\|
 \leq C\rho^{-n}v\,v^{-2}
 =C(2/\rho)^n\longrightarrow0.
\]
The left side is independent of $n$, so $\mu A_{2+i}=\mu$.
Composition with $A_{-2+i}$ gives the negative period as well.

\paragraph{Horizontal and scale averages.}
Define, by integration on continuous test functions,
\[
 \bar\mu=\frac12\int_0^2\mu A_{x+i}\dd x.
\]
For each continuous test, the integrand is continuous in $x$, by joint
continuity of the normalized transformations and their weights.  Compactness
of the parameter interval makes this a finite positive functional, and its
mass is strictly positive because the weights $|Q_F(x+i)|^2$ are positive.
The period just proved shows that
\begin{equation}
 \bar\mu A_{t+i}=\bar\mu\qquad(t\in\mathbb R).
 \label{eq:horizontal-average}
\end{equation}

We next compute the change under a factor of two in height.  Using
\eqref{eq:transfer-eigenmeasure} before expanding $L$, we obtain
\begin{align*}
 \rho\bar\mu
 &=\frac12\int_0^2\mu L A_{x+i}\dd x\\
 &=\frac14\int_0^2
      \mu\bigl(A_{-1/2+x/2+i/2}+A_{1/2+x/2+i/2}\bigr)\dd x\\
 &=\frac12\int_{-1/2}^{3/2}\mu A_{s+i/2}\dd s.
\end{align*}
Here the substitutions $s=\pm1/2+x/2$ give two adjacent intervals.
The last integrand has period two, since
\[
 \mu A_{s+2+i/2}=\mu A_{2+i}A_{s+i/2}
                =\mu A_{s+i/2}.
\]
Moving its integration interval to $[0,2]$ yields
\begin{equation}
 \rho\bar\mu=\bar\mu A_{i/2}.
 \label{eq:half-scale-average}
\end{equation}
In particular
\[
 \bar\mu A_{2iv}=\rho^{-1}\bar\mu A_{iv},
 \qquad v>0.
\]
Since $2^\beta=\rho$, the functional
\[
 H(v)=v^\beta\bar\mu A_{iv}
\]
satisfies $H(2v)=H(v)$.  Moreover, each $\bar\mu A_{iv}$ remains
horizontally invariant: the affine product gives
\[
 A_{iv}A_{t+i}=A_{vt+i}A_{iv},
\]
and \eqref{eq:horizontal-average} applies to the first operator on the
right.

We have obtained all horizontal translations and a dyadic scaling relation.
To obtain every positive scale, set
\[
 \nu=\int_1^2 H(v)\,\frac{dv}{v}.
\]
As before, this is a finite positive functional of strictly positive mass.
For every $y>0$, substitution $u=vy$ gives
\begin{align*}
 \nu A_{iy}
 &=y^{-\beta}\int_y^{2y}H(u)\,\frac{du}{u}
  =y^{-\beta}\nu.
\end{align*}
The last equality follows because $t\mapsto H(e^t)$ has period $\log2$;
the integral over one full period is independent of its starting point.
Horizontal invariance also persists under this average.  Normalize to the
probability $P=\nu/\nu(\mathbf 1)$.  Since
$A_{x+iy}=A_{x+i}A_{iy}$, we conclude that
\[
 PA_{x+iy}=y^{-\beta}P
\]
on continuous tests.

\paragraph{Borel tests.}
For each fixed $z=x+iy$, the expression
\[
 E\longmapsto\int_{\KK}|Q_F(z)|^2
                \mathbf 1_E(T_zF)\dd P(F)
\]
defines a finite Borel measure on $\KK$.  Indeed, its weight and
transformation are continuous in $F$, and the weight is bounded on $\KK$.
The identity on continuous functions identifies this measure with
$y^{-\beta}P$.  Equality on nonnegative simple functions, followed by
monotone convergence, proves \eqref{eq:weighted-covariance} in its stated
generality.
\end{proof}

The covariance identity also applies to integrable signed tests, by their
positive and negative parts.  When a nonnegative test depends in a jointly
Borel way on an additional parameter, we may apply the identity for each
fixed parameter and then integrate by Tonelli's theorem.  These uses require
only the equality of weighted measures just proved.

\begin{lemma}[Moments of the logarithmic derivative]
\label{lem:law-moments}
Assume $\rho>2$, and let $P$ and $\beta$ be as in
Lemma~\ref{lem:weighted-law}.  For $z=x+iy\in\HH$, define
\begin{equation}
 q_F(z)=iy\frac{Q'_F(z)}{Q_F(z)}
       =-iy\frac{F''(z)}{F'(z)}
       =a_F(z)+ib_F(z),
 \label{eq:logarithmic-derivative-variable}
\end{equation}
where $a_F,b_F$ are real.  The variables $q_F(i)$ are bounded on $\KK$, and
\begin{equation}
 \EE a_F(i)=-\frac\beta2,
 \qquad \EE b_F(i)=0,
 \qquad \EE|q_F(i)|^2=\frac{\beta(\beta+1)}4.
 \label{eq:law-moments}
\end{equation}
They also satisfy
\begin{equation}
 q_{T_zF}(w)=q_F(z\cdot w),\qquad z,w\in\HH.
 \label{eq:q-affine-covariance}
\end{equation}
\end{lemma}

\begin{proof}
Taking $\Phi=\mathbf 1$ in \eqref{eq:weighted-covariance} gives
\begin{equation}
 \EE|Q_F(x+iy)|^2=y^{-\beta}.
 \label{eq:inverse-derivative-mean}
\end{equation}
On a compact neighborhood of any interior point, the functions
$Q_F,Q'_F,Q''_F$ are uniformly bounded over $\KK$, by
Lemma~\ref{lem:compact-class}.  These bounds control the first two real
derivatives of $|Q_F|^2$, and hence justify differentiation under expectation
in \eqref{eq:inverse-derivative-mean}.

Holomorphicity and \eqref{eq:logarithmic-derivative-variable} give the
pointwise identities
\[
 \partial_y|Q_F|^2=\frac{2a_F}{y}|Q_F|^2,
 \qquad
 \partial_x|Q_F|^2=\frac{2b_F}{y}|Q_F|^2,
 \qquad
 \Delta|Q_F|^2=4|Q'_F|^2.
\]
At $i$ we have $Q_F(i)=1$, so the $y$ derivative, the $x$ derivative, and
the Euclidean Laplacian of \eqref{eq:inverse-derivative-mean} give,
respectively,
\[
 2\EE a_F(i)=-\beta,
 \qquad 2\EE b_F(i)=0,
 \qquad 4\EE|q_F(i)|^2=\beta(\beta+1).
\]
This proves \eqref{eq:law-moments}.  Boundedness of $q_F(i)$ follows from
the same compact-family bounds and $Q_F(i)=1$.

Finally, differentiating
$Q_{T_zF}(w)=Q_F(z\cdot w)/Q_F(z)$ with respect to $w$ introduces the factor
$y=\operatorname{Im}z$.  Since
$\operatorname{Im}(z\cdot w)=y\operatorname{Im}w$, formula
\eqref{eq:q-affine-covariance} follows directly from the definition of $q$.
\end{proof}

The law has now supplied affine covariance and explicit moments.
We next use it to control the behavior of typical maps near the boundary
and at infinity; the moments will enter the later geometric comparison.
